\section{Data Feed - AI Evaluierung}
Siehe \href{sec:KI}{KI Nutzung} für details, Link zum\href{https://gemini.google.com/share/5f8f9fe4f518}{Prompt}\\

\noindent
Hallo! Als potenzieller Nutzer finde ich die Grundidee von "Data Feed" extrem spannend. Die Vision, eine lehrreiche Alternative zum klassischen "Doomscrolling" zu schaffen, trifft genau den Zeitgeist. Hier ist mein Feedback zu eurem Konzept und Prototyp:
\\

\noindent
Was mir sehr gut gefällt:
\\

\noindent
\emph{Visuelle Klarheit:} Der Prototyp wirkt wunderbar aufgeräumt. Die Idee, komplexe Daten in kleine, snackbare "Posts" herunterzubrechen, ist super übersichtlich umgesetzt.
\\

\noindent
\emph{Barrierefreiheit:} Die starke Fokussierung auf Inklusion (anpassbare Kontraste, Berücksichtigung von Farbenblindheit) direkt im Menü ist ein riesiger Pluspunkt!
\\

\noindent
\emph{Navigation:} Das Themen-Karussell am unteren Rand wirkt intuitiv. Auch die geplante "Swipe-down"-Geste zum Schließen der Detailansicht ist modern und nutzerfreundlich. Die Einbindung bekannter APIs (wie Google Maps für die Baumkarte) senkt die Einstiegshürde enorm.
\\

\noindent
Wo ich noch kritisch bin:
\\

\noindent
\emph{Der "Endlos"-Widerspruch:} Ihr beschreibt die App als entspannte Alternative zum Doomscrolling, nutzt aber als Kernmechanik einen endlosen Feed. Das widerspricht sich etwas. Endloses Scrollen verleitet oft zum gedankenlosen Wischen, selbst bei guten Inhalten. Vielleicht wären tägliche, limitierte "Daten-Häppchen" oder ein klarer Endpunkt entspannender?
\\

\noindent
\emph{Inhaltsauswahl:} Die App steht und fällt mit der Kuration. Die "Baumkarte Wien" ist ein tolles, greifbares Beispiel. "Jagdvergabe durch die Jahre" (aus dem Prototyp) wirkt hingegen auf den ersten Blick sehr trocken für einen Casual-User. Die begleitenden Erklärtexte müssen wirklich extrem gut und simpel geschrieben sein, damit es unterhaltsam bleibt.
\\

\noindent
\emph{Fazit:} Ein visuell starkes und durchdachtes Konzept mit einem großartigen Ziel. Wenn ihr sicherstellt, dass die Feed-Mechanik nicht doch wieder Stress erzeugt und die Daten wirklich alltagstauglich und spannend aufbereitet sind, würde ich diese App definitiv herunterladen!

